\begin{theorem}[Commuting closure labels and centralizers]%
    \label{thm:peps_pair_centralizer_count}
    \lean{TNLean.PEPS.card_commutingPairConjugacyClass_eq_sum_card_centralizer_of_representatives,
        TNLean.PEPS.card_commutingPairConjugacyClass_eq_sum_card_centralizer,
        TNLean.PEPS.commutingPairConjugacyClassEquivSigmaCentralizerOfRepresentatives,
        TNLean.PEPS.commutingPairConjugacyClassEquivSigmaCentralizer,
        TNLean.PEPS.centralizerConjClassEquivCommutingPairFiber}
    \leanok
    \uses{def:peps_pair_conjugacy}
    Let $G$ be a group and choose a representative $g_C$ of each
    conjugacy class $C$. Write $\operatorname{Cl}(H)$ for the set of
    conjugacy classes of a group $H$. The simultaneous-conjugacy classes
    of commuting pairs are in bijection with
    $\coprod_{C\in\operatorname{Cl}(G)}\operatorname{Cl}(Z_G(g_C))$.
    If $G$ is finite, their number is therefore
    \begin{align}
        K=\sum_{C\in\operatorname{Cl}(G)}
        \bigl|\operatorname{Cl}(Z_G(g_C))\bigr|.
        \label{eq:peps_pair_class_count}
    \end{align}
    This is the group-theoretic classification following
    \cite[Theorem~5.9]{Schuch2010PEPS}.
\end{theorem}

\begin{proof}\leanok
    \uses{thm:peps_pair_conjugacy_centralizer}
    Assign to the class of $(g,h)$ the conjugacy class of $g$.
    If $xgx^{-1}=g_C$, then $xhx^{-1}\in Z_G(g_C)$ and
    $(g,h)$ is simultaneously conjugate to $(g_C,xhx^{-1})$.
    For a fixed first element $g_C$, two such pairs have the same
    simultaneous-conjugacy class exactly when their second elements
    are conjugate in $Z_G(g_C)$, by
    Theorem~\ref{thm:peps_pair_conjugacy_centralizer}.
    This proves the stated bijection. Counting its disjoint components
    gives (\ref{eq:peps_pair_class_count}).
\end{proof}

\begin{theorem}[Dimension of the commuting torus closure span]%
    \label{thm:peps_g_injective_closure_dimension}
    \lean{TNLean.PEPS.IsGInjective.finrank_span_torusGClosure_commuting_eq_sum_card_centralizer,
        TNLean.PEPS.IsGInjective.finrank_span_torusGClosureClass_commuting_eq_sum_card_centralizer,
        TNLean.PEPS.IsGInjective.finrank_span_torusGClosureClass_commuting_of_isSemiRegular,
        TNLean.PEPS.range_torusGClosure_commuting_eq_range_torusGClosureClass}
    \leanok
    \uses{def:peps_semiregular, def:peps_g_injective,
        def:peps_torus_g_closure, def:peps_pair_conjugacy}
    Let $U$ be a semi-regular representation of a finite group $G$,
    and let $A$ be $G$-injective for its four-leg virtual representation.
    On a nonempty rectangular torus, let $\mathcal V_A$ be the span
    of the vectors $\ket{\mathcal M(A\mid C)}$ indexed by commuting
    simultaneous-conjugacy classes $C$. Then
    \begin{align}
        \dim\mathcal V_A
        =\sum_{C\in\operatorname{Cl}(G)}
        \bigl|\operatorname{Cl}(Z_G(g_C))\bigr|.
    \end{align}
    This is the dimension of the commuting closure span appearing in
    \cite[Theorem~5.9]{Schuch2010PEPS}.
\end{theorem}

\begin{proof}\leanok
    \uses{thm:peps_semiregular_g_injective_closure_independence,
        thm:peps_pair_centralizer_count}
    The class-labelled vectors are linearly independent by
    Theorem~\ref{thm:peps_semiregular_g_injective_closure_independence}.
    They form a basis of their span, whose dimension is therefore
    the number of commuting pair classes. Apply
    Theorem~\ref{thm:peps_pair_centralizer_count}.
\end{proof}
